


\section{Problem Formulation}
\label{subsec:problem_formulation}
\subsection{QHD: hyperdimensional Q-learning for continuous-state control}
We now describe how hyperdimensional representations are used to approximate action-value functions in continuous-state, discrete-action control. Agent $i$ selects an HDC encoder
\begin{equation}
\Phi_i:\mathcal{S}\to\mathbb{K}^{D_i},
\qquad \mathbb{K}\in\{\mathbb{R},\mathbb{C}\},
\label{eq:phi_i}
\end{equation}
which maps continuous states $s\in\mathbb{R}^{d}$ to hypervectors. For each action $a\in\mathcal{A}$, the agent maintains a parameter hypervector $\mathbf{w}_{i,a}\in\mathbb{K}^{D_i}$ and defines the Q-function approximation
\begin{equation}
Q_i(s,a) \;\triangleq\; \Re\!\left(\langle \Phi_i(s),\,\mathbf{w}_{i,a}\rangle\right).
\label{eq:qhd_q}
\end{equation}
Stacking parameters across actions yields $W_i [\mathbf{w}_{i,a}]_{a\in\mathcal{A}}\in\mathbb{K}^{D_i\times|\mathcal{A}|}$ and $Q_i(s,\cdot)=\Re(\Phi_i(s)^{\mathrm{H}} W_i)$. Although the encoder $\Phi_i(\cdot)$ may be highly nonlinear in $s$,
the model is \emph{linear in parameters}. This linearity is the key structural property we will exploit for federated learning.

Given a replayed transition $(s,a,r,s')$ and a stabilized target estimate (e.g., Double Q-learning with a delayed target copy $Q_i^{-}$), define
\begin{equation}
y = r + \gamma\, Q_i^{-}\!\left(s',\,\arg\max_{a'\in\mathcal{A}} Q_i(s',a')\right).
\label{eq:target}
\end{equation}
QHD performs a lightweight rank-one update
\begin{equation}
\mathbf{w}_{i,a} \leftarrow \mathbf{w}_{i,a} + \eta\,\big(y - Q_i(s,a)\big)\,\Phi_i(s),
\label{eq:qhd_update}
\end{equation}
where $\eta>0$ is a step size...

{\color{red} explain the updates}





% \subsection{HDC Q-function Representation}
% \label{subsec:hdc_q_representation}

% Hyperdimensional computing (HDC) encodes symbols using high-dimensional randomized
% hypervectors with near-orthogonality properties, enabling rich compositional structure while
% preserving simple arithmetic operations. In our setting, an HDC encoder
% \(
%     \Phi: \mathcal{S}\times\mathcal{A}\to\mathbb{C}^D
% \)
% maps each state--action pair $x=(s,a)$ to a $D$-dimensional hypervector $\Phi(x)$.
% The encoder is sampled once at initialization and shared across all agents.
% To ensure numerical stability and compatibility with our finite-time analysis, we assume
% \(
%     \|\Phi(s,a)\|_2 \le 1
%     \,\text{for all $(s,a)$}.
% \)
% Each agent $i$ maintains a parameter hypervector $\mathbf{M}^{(i)}\in\mathbb{C}^D$ and evaluates its
% action--value function as
% \begin{equation}
%     Q_{\mathbf{M}}(s,a)
%     \;=\;
%     \langle \Phi(s,a),\, \mathbf{M} \rangle,
%     \label{eq:hdc_q_function}
% \end{equation}
% where $\langle\cdot,\cdot\rangle$ is the standard inner product on $\mathbb{C}^D$.
% This parameterization exhibits three key structural properties:

% \textit{1)Linearity in parameters:}
%     The map $M \mapsto Q_M$ is linear, so the HDC semi-gradient update 
%     reduces to a rank-one adjustment with $O(D)$ cost, matching the 
%     computational footprint of linear federated TD methods.

% \textit{2)Nonlinearity in inputs:}
%     The encoder $\Phi(s,a)$ is generally nonlinear in $(s,a)$, enabling 
%     $Q_M$ to capture dependencies that cannot be expressed by linear 
%     features. As a result, the HDC representation behaves like a 
%     high-capacity function class while retaining computational simplicity.

% \textit{3)Random-feature kernel structure:}
%     The induced kernel $k_D(x,x')=\langle \Phi(x),\Phi(x')\rangle$ 
%     concentrates to a smooth limiting kernel $k_\ast$ as $D\to\infty$. 
%     Thus, $Q_M$ can be viewed as a finite-dimensional approximation to an 
%     RKHS value function.




\subsection{Federated value learning as a global objective}
We consider $N$ agents learning concurrently, each interacting with its own environment and collecting a private replay buffer $\mathcal{D}_i$. ... 


Federated reinforcement learning seeks to aggregate local value knowledge into a global representation without centralizing data, data sharing is impossible due to privacy, communication limitations, latency, etc. 




\paragraph{Federation as function-space averaging (synchronization operator).}
Each client $i$ learns a local action-value model $Q_i$ from private off-policy replay using QHD-style updates. The server does not access client data or optimize client Bellman losses; instead, it periodically performs a \emph{synchronization} step that aggregates the learned client models into a global teacher in \emph{function space}. Let $\mu$ denote an agreement/query distribution over state--action pairs (instantiated in our method by a small shared anchor set), and let $\pi_i\ge 0$ be federation weights with $\sum_{i=1}^N \pi_i=1$ (e.g., proportional to client data volume).
Given fixed client models $\{Q_i\}_{i=1}^N$, the server defines the global teacher $Q^{\mathrm{glob}}$ as the unique minimizer of
\begin{equation}
Q^{\mathrm{glob}} \;\in\; \arg\min_{Q} \sum_{i=1}^N \pi_i\, E_{(s,a)\sim\mu}\!\left[\big(Q_i(s,a)-Q(s,a)\big)^2\right].
\label{eq:fed_sync}
\end{equation}
This objective is to combine the clients' knowledge in a shared output space rather than in an (often incompatible) parameter space.  The minimizer is the weighted pointwise average (in $\mu$-expectation),
\begin{equation}
Q^{\mathrm{glob}}(s,a) \;=\; \sum_{i=1}^N \pi_i\,Q_i(s,a) \qquad (\mu\text{-a.e.}).
\label{eq:qglob_avg}
\end{equation}

{\color{red} introduce the general problem of function space Q Learning, which can be implied to tabular or continuous state-action setting, to linear or non-linear paprameter averaging in reinforcement learning methond. claim the obstacless}


For TABLUAR, THE solution on top is like hsaring their Q-table, which leads to the exact function-sapce soltuions knowln asn Federated Q-Learning.  FEDERATED AVERAGING HOLDS THE EXACT FUNCTION SPACE, SIMILAR TO A LINEAR FORM LIKE QHD. BUT, FOR NONLINEAR, THEY RELY ON THE HEURISTIC PARAMETER AVERAGING. 



Existing federated RL systems face two central obstacles:{(1) No closed-form, fast synchronization under nonlinear models.} For nonlinear approximators, federated aggregation generally cannot be computed in closed form; client-side assimilation requires iterative optimization, which is expensive and unstable under moving RL targets. {(2) Model heterogeneity breaks parameter aggregation.} Classical federated methods assume a shared parameterization. If clients use different encoders, different architectures, or different parameter dimensions $D_i$, parameter averaging is undefined or semantically meaningless. Such heterogeneity is unavoidable in real autonomous systems.



\section{FedQHD: Function-Space Federation with Closed-Form Translation}
\label{sec:methodology}
The federation objective in Eq.~\eqref{eq:fed_sync} defines a synchronization operator that aggregates client knowledge in \emph{function space}, independently of how local value functions are parameterized or learned. In this section, we show that QHD provides a uniquely favorable structure for realizing this operator efficiently and exactly.

The key enabling property is that, although QHD induces expressive nonlinear approximation over continuous state spaces through the encoder $\Phi_i(\cdot)$, the resulting Q-functions are \emph{linear in parameters}. This linearity allows the federation objective to be solved in closed form when representations are shared, and to be extended to heterogeneous encoders via closed-form translation operators. We first analyze the homogeneous-encoder regime, where the federation operator admits an exact solution that coincides with parameter averaging. This case clarifies the relationship between FedQHD and FedAvg and motivates the general heterogeneous construction that follows.

\subsection{Fed-QHD with Homogeneous Encoders}
\label{sec:fedqhd_homogeneous}

\begin{equation}
Q^{\mathrm{glob}}(s,a) = \langle \Phi(s), w^{\mathrm{glob}}_a \rangle,
\qquad
W^{\mathrm{glob}} \in \mathbb{K}^{D \times |\mathcal{A}|}.
\label{eq:qhd_global}
\end{equation}
Substituting \eqref{eq:qhd_local} and \eqref{eq:qhd_global} into \eqref{eq:fed_sync_recall} yields the quadratic objective
\begin{equation}
\min_{W^{\mathrm{glob}}}\sum_{i=1}^N \pi_i E_{(s,a)\sim\mu} \big[ \langle \Phi(s), (W_i - W^{\mathrm{glob}})_{:,a} \rangle^2\big].
\label{eq:fed_quad}
\end{equation}
Let $\Sigma_\mu = \E_{(s,a)\sim\mu}[\Phi(s)\Phi(s)^\mathrm{H}]$ denote the
feature covariance induced by the agreement distribution.
Then \eqref{eq:fed_quad} can be written as
\begin{equation}
\min_{W^{\mathrm{glob}}} \sum_{i=1}^N \pi_i \sum_{a\in\mathcal{A}} (W_{i,a}-W^{\mathrm{glob}}_a)^\mathrm{H} \Sigma_\mu (W_{i,a}-W^{\mathrm{glob}}_a),
\label{eq:fed_quad_cov}
\end{equation}
which is a strictly convex quadratic problem whenever $\Sigma_\mu \succ 0$, with solution as:
\begin{equation}
W^{\mathrm{glob}} = \sum_{i=1}^N \pi_i W_i.
\label{eq:exact_avg}
\end{equation}

{\color{red} remind that $\Sigma_\mu$ doesn't explicitly appear in solution but involves the raw data sharing with homo encoders}


\subsection{Fed-QHD with Heterogeneous Encoders}
\label{sec:fedqhd_heterogeneous}
Consider a case where each client may use a different hyperdimensional encoder, different kernel hyperparameters, and even a different feature dimension.  In this setting, parameter averaging is undefined, but function-space consensus remains meaningful. Client $i$ uses an encoder \(\Phi_i:\mathcal{S}\to\mathbb{K}^{D_i}, \mathbb{K}\in\{\mathbb{R},\mathbb{C}\},\)
where both the dimension $D_i$ and encoder design (random seed, binding/permutation rules, kernel scale, quantization)
may differ across clients. 

To define federation operationally, we instantiate the agreement distribution $\mu$ by a finite \emph{server-defined anchor set}. Let $\mathcal{S}_{\mathrm{ref}}=\{s_1,\dots,s_m\}$ denote a fixed set of reference states generated independently of all clients. In many autonomy domains, the server can obtain $\mathcal{S}_{\mathrm{ref}}$ from a simulator/digital twin by executing a random or exploratory policy and logging visited states without rewards, actions, or trajectories.

Each client computes its anchor predictions $Q^{\mathrm{ref}}_i[j,a] \;\defeq\; Q_i(s_j,a),\, Q^{\mathrm{ref}}_i\in\mathbb{R}^{m\times|\mathcal{A}|},$
and the server forms the function-space consensus teacher
$Q^{\mathrm{glob}}_{\mathrm{ref}}=\sum_{i=1}^N \pi_i\, Q^{\mathrm{ref}}_i, \sum_{i=1}^N \pi_i=1,\;\pi_i\ge 0$,
which is the unique minimizer of the anchor consensus objective
$\min_{Q\in\mathbb{R}^{m\times|\mathcal{A}|}} \sum_i \pi_i \|Q^{\mathrm{ref}}_i - Q\|_F^2$.


DEFINE AN ANCHOR FOR CLIENT I

For each client $i$, define the \emph{anchor feature matrix} by stacking encoded anchor states:
\begin{equation}
X_i \;=\; 
\begin{bmatrix}
\Phi_i(s_1) \\
\Phi_i(s_2) \\
\vdots \\
\Phi_i(s_m)
\end{bmatrix}
\in \mathbb{K}^{m \times D_i}.
\label{eq:anchor_matrix_i}
\end{equation}
The local Q-function evaluations on the anchor set can then be written compactly as
\begin{equation}
Q^{\mathrm{ref}}_i = \Re(X_i W_i),
\qquad W_i = [\mathbf{w}_{i,a}]_{a\in\mathcal{A}} \in \mathbb{K}^{D_i \times |\mathcal{A}|}.
\label{eq:anchor_q_matrix_form}
\end{equation}

THEN WRITE THE REGULARIZED NONLINEAR EQUATIONS

We seek a client-side parameter matrix $W^{\mathrm{glob}}_i\in\mathbb{K}^{D_i\times|\mathcal{A}|}$ that enables client $i$ to reproduce the global teacher $Q^{\mathrm{glob}}_{\mathrm{ref}}$ as closely as possible using its own encoder $\Phi_i$. This leads to the \emph{regularized least-squares compilation problem}:
\begin{equation}
W^{\mathrm{glob}}_i \;=\; \arg\min_{W \in \mathbb{K}^{D_i \times |\mathcal{A}|}} 
\left\|X_i W - Q^{\mathrm{glob}}_{\mathrm{ref}}\right\|_F^2 
+ \lambda \|W\|_F^2,
\label{eq:regularized_compilation}
\end{equation}
where $\lambda \ge 0$ is a regularization parameter that controls model complexity and stabilizes the solution when the anchor set is small ($m < D_i$) or when $X_i$ is ill-conditioned.


AND THE CONVEX SOLUTIUON FOR IT. 

Problem~\eqref{eq:regularized_compilation} is a convex quadratic optimization problem in $W$ with a unique global minimizer. Taking the gradient with respect to $W$ and setting it to zero yields the regularized normal equations:
\begin{equation}
\left(X_i^\mathrm{H} X_i + \lambda I_{D_i}\right) W^{\mathrm{glob}}_i = X_i^\mathrm{H} Q^{\mathrm{glob}}_{\mathrm{ref}}.
\label{eq:normal_equations}
\end{equation}
When $\lambda > 0$ or when $X_i^\mathrm{H} X_i$ is full rank, the system~\eqref{eq:normal_equations} is invertible, yielding the explicit solution:
\begin{equation}
W^{\mathrm{glob}}_i = \left(X_i^\mathrm{H} X_i + \lambda I_{D_i}\right)^{-1} X_i^\mathrm{H} Q^{\mathrm{glob}}_{\mathrm{ref}}.
\label{eq:ridge_solution}
\end{equation}
In the unregularized limit $\lambda \to 0$, this reduces to the Moore-Penrose pseudoinverse solution:
\begin{equation}
W^{\mathrm{glob}}_i = X_i^\dagger Q^{\mathrm{glob}}_{\mathrm{ref}},
\qquad X_i^\dagger = (X_i^\mathrm{H} X_i)^{-1} X_i^\mathrm{H},
\label{eq:pseudoinverse_solution}
\end{equation}
provided $\mathrm{rank}(X_i) = \min(m, D_i)$.



SEE IF YOU CAN CONNECT HERE WITH ONE EQUATION TO GRAM MATRIX


PROCESS: CLIENT SHARE THEIR DATA, THE SERVER COMPUTES THE Q GLOBAL, AND RETURN THE .... 



% While the shared encoder approach provides efficient aggregation through direct parameter
% averaging, practical federated deployments often involve agents with heterogeneous
% computational capabilities or environmental complexities. To support such scenarios, we
% extend FedQHD to accommodate \textit{agent-specific encoders} $\{\Phi^{(i)}\}_{i=1}^N$
% while preserving the ability to aggregate learned knowledge.

% \subsubsection{Motivation for Heterogeneous Encoders}

% Different agents may benefit from different encoder architectures for several reasons:
% \begin{itemize}
%     \item \textbf{Computational constraints:} Edge devices may require lower-dimensional
%     encoders than datacenter agents.
%     \item \textbf{Environmental complexity:} Agents operating in complex environments may
%     need richer feature representations.
%     \item \textbf{Local adaptation:} Fine-tuning encoders to local state distributions can
%     improve sample efficiency.
% \end{itemize}

% However, with heterogeneous encoders, direct parameter averaging is no longer valid since
% $\mathbf{M}^{(i)}$ lives in the feature space induced by $\Phi^{(i)}$. We address this
% challenge through \textit{regression-based aggregation}.









\section{Projection Residual as a Structural Heterogeneity Measure}
\label{sec:projection_residual}

We formalize the least-squares compilation residual induced by heterogeneous client representations and show that it admits an exact operator form with a clear geometric and optimization-theoretic interpretation. This residual isolates irreducible representation mismatch across clients and provides a principled diagnostic for capacity-limited transfer in heterogeneous federated reinforcement learning.


\begin{proposition}[Irreducible anchor mismatch]
\label{prop:irreducible}
Assume $\lambda=0$ and let $P_i$ denote the orthogonal projector onto
$\mathrm{col}(X_i)\subseteq\mathbb{K}^m$. Then
\[
\min_{W\in\mathbb{K}^{D_i\times|A|}}\|X_iW-Q^{\mathrm{glob}}_{\mathrm{ref}}\|_F
\;=\;\|(I-P_i)Q^{\mathrm{glob}}_{\mathrm{ref}}\|_F.
\]
In particular, $\|(I-P_i)Q^{\mathrm{glob}}_{\mathrm{ref}}\|_F=0$ if and only if
$Q^{\mathrm{glob}}_{\mathrm{ref}}\in\mathrm{col}(X_i)^{|A|}$, i.e., the teacher is
exactly representable by client $i$ on the anchor set.
\end{proposition}


\begin{proposition}[Residual amplification via subspace misalignment]
\label{prop:heterogeneity}
Assume $\lambda=0$ and let $P_i$ denote the orthogonal projector onto
$\mathrm{col}(X_i)$. Suppose the anchor teacher is formed as a convex
combination
\(
Q^{\mathrm{glob}}_{\mathrm{ref}}
=
\sum_{j=1}^N \pi_j Q^{(j)}_{\mathrm{ref}},
\qquad
Q^{(j)}_{\mathrm{ref}}\in\mathrm{col}(X_j),
\quad
\sum_j\pi_j=1.
\)
Then the irreducible mismatch for client $i$ satisfies
\[
R_{i,0}= (I-P_i)Q^{\mathrm{glob}}_{\mathrm{ref}} = \sum_{j\neq i}\pi_j (I-P_i)Q^{(j)}_{\mathrm{ref}}, \text{ with the bound } \|R_{i,0}\|_F \;\le\; \sum_{j\neq i}\pi_j \sin(\theta_{ij}) \,\|Q^{(j)}_{\mathrm{ref}}\|_F,
\]
where $\theta_{ij}$ is the largest principal angle between the subspaces
$\mathrm{col}(X_i)$ and $\mathrm{col}(X_j)$.
\end{proposition}

% \subsection{Federated Learning Protocol}
% \label{subsec:federated_protocol}

% With a shared encoder $\Phi$ across all agents, the linearity of HDC Q-functions in the
% parameter space enables a simple and efficient federated aggregation strategy: direct
% parameter averaging. This approach leverages the parameter-function equivalence property
% (Theorem~\ref{thm:parameter_function_equivalence}), which guarantees that averaging
% parameter vectors is equivalent to averaging Q-functions pointwise.

% \subsection{Direct Parameter Averaging}
% \label{subsec:federated_averaging}
% FedQHD operates in synchronous communication rounds indexed by $k=0,1,2,\dots$. Each
% round consists of three phases with $K$ local update steps, so the global iteration count and
% round index are related by $t = kK + u$, where $u \in \{0,\dots,K-1\}$ denotes
% the local step within round $k$.

% \subsubsection{Broadcast}
% The server maintains a global hypervector $\mathbf{M}_k$ and broadcasts it along with the
% shared encoder $\Phi$ to all clients. Each client initializes its local parameter to $\mathbf{M}_k$.

% \subsubsection{Local Updates}
% Agent $i$ executes $K$ on-policy HDC TD updates using the shared encoder $\Phi$,
% generating
% \begin{equation}
% \begin{split}
%     \mathbf{M}^{(i)}_{k,0}=\mathbf{M}_k,
%     \quad
%      \mathbf{M}^{(i)}_{k,u+1}
%     =
%     \mathbf{M}^{(i)}_{k,u}
%     + \alpha_{k,u}\, \delta_t^{(i)} \Phi(s_t^{(i)}, a_t^{(i)})
% \end{split}
%     \label{eq:local_update}
% \end{equation}
% for $u=0,\dots,K-1$, where $\delta_t^{(i)}$ is the TD error computed using the local Q-function
% $Q_{\mathbf{M}^{(i)}_{k,u}}(s,a) = \langle \Phi(s,a), \mathbf{M}^{(i)}_{k,u} \rangle$.

% \subsubsection{Upload and Aggregation}
% The server receives the local parameter vectors $\{\mathbf{M}^{(i)}_{k,K}\}_{i=1}^N$ and performs
% direct federated averaging:
% \begin{equation}
%     \mathbf{M}_{k+1}
%     =
%     \frac{1}{N}\sum_{i=1}^{N}
%     \mathbf{M}^{(i)}_{k,K}.
%     \label{eq:federated_averaging}
% \end{equation}
% By Theorem~\ref{thm:parameter_function_equivalence}, this parameter averaging is equivalent
% to averaging the Q-functions, i.e., $Q_{\mathbf{M}_{k+1}}(s,a) = \frac{1}{N}\sum_{i=1}^N Q_{\mathbf{M}^{(i)}_{k,K}}(s,a)$
% for all $(s,a)$. The averaged model $\mathbf{M}_{k+1}$ is then broadcast to all clients to
% initiate round $k+1$.
\subsection{FedQHD Algorithm}
\label{subsec:algorithm}

% We present the complete FedQHD algorithm with shared encoder and direct parameter averaging.
% This approach leverages the parameter-function equivalence property to enable efficient
% federated aggregation with minimal communication overhead ($O(D)$ parameter uploads per round).

\begin{algorithm}[t]
\caption{FedQHD: Shared Encoder Protocol}
\label{alg:fedqhd_shared}
\begin{algorithmic}[1]
\STATE \textbf{Input:} Rounds $T$, local steps $K$, agents $N$, learning rates $\{\alpha_{k,u}\}$
\STATE \textbf{Server Initialize:}
\STATE \quad Shared encoder $\Phi: \mathcal{S}\times\mathcal{A}\to\mathbb{C}^D$
\STATE \quad Global model $\mathbf{M}_0 = \mathbf{0} \in \mathbb{C}^D$
\STATE
\FOR{round $k = 0, 1, \ldots, T-1$}
    \STATE \textbf{[Server Broadcast]}
    \STATE \quad Send $\mathbf{M}_k$ to all agents (encoder $\Phi$ sent once at initialization)
    \STATE
    \STATE \textbf{[Client Local Learning]} (in parallel for all agents $i = 1, \ldots, N$)
    \STATE \quad Initialize: $\mathbf{M}^{(i)}_{k,0} = \mathbf{M}_k$
    \FOR{local step $u = 0, 1, \ldots, K-1$}
        \STATE \quad Observe state $s_t^{(i)}$, select action $a_t^{(i)}$ via $\epsilon$-greedy on $Q_{\mathbf{M}^{(i)}_{k,u}}$
        \STATE \quad Execute $a_t^{(i)}$, observe reward $r_t^{(i)}$ and next state $s_{t+1}^{(i)}$
        \STATE \quad Select next action $a_{t+1}^{(i)}$ using current policy (SARSA)
        \STATE \quad Encode: $\phi_t = \Phi(s_t^{(i)}, a_t^{(i)})$, $\phi_{t+1} = \Phi(s_{t+1}^{(i)}, a_{t+1}^{(i)})$
        \STATE \quad TD error: $\delta_t^{(i)} = r_t^{(i)} + \gamma \langle \phi_{t+1}, \mathbf{M}^{(i)}_{k,u} \rangle - \langle \phi_t, \mathbf{M}^{(i)}_{k,u} \rangle$
        \STATE \quad \textbf{Update:} $\mathbf{M}^{(i)}_{k,u+1} = \mathbf{M}^{(i)}_{k,u} + \alpha_{k,u} \delta_t^{(i)} \phi_t$
    \ENDFOR
    \STATE \quad \textbf{Upload:} Send $\mathbf{M}^{(i)}_{k,K}$ to server
    \STATE
    \STATE \textbf{[Server Aggregation]}
    \STATE \quad Receive parameter vectors $\{\mathbf{M}^{(i)}_{k,K}\}_{i=1}^{N}$ from all agents
    \STATE \quad \textbf{Average parameters:} $\mathbf{M}_{k+1} = \frac{1}{N}\sum_{i=1}^{N} \mathbf{M}^{(i)}_{k,K}$
    \STATE \quad \textbf{(Optional)} Project: $\mathbf{M}_{k+1} \leftarrow \text{Proj}_{\mathcal{B}(G)}(\mathbf{M}_{k+1})$
\ENDFOR
\STATE
\STATE \textbf{Output:} Global Q-function $Q_{\mathbf{M}_T}(s,a) = \langle \Phi(s,a), \mathbf{M}_T \rangle$
\end{algorithmic}
\end{algorithm}

% \subsection{Theoretical Property: Parameter-Function Equivalence}
% \label{subsec:theory}

% \begin{theorem}[Parameter-Function Equivalence]
% \label{thm:parameter_function_equivalence}
% Let $\Phi: \mathcal{S}\times\mathcal{A} \to \mathbb{C}^D$ be a shared encoder fixed across all agents, and let $\{\mathbf{M}^{(i)}\}_{i=1}^N$ be local model hypervectors. Define the averaged hypervector as
% \begin{equation}
%     \bar{\mathbf{M}} = \frac{1}{N}\sum_{i=1}^N \mathbf{M}^{(i)}.
% \end{equation}
% Then, for all state-action pairs $(s,a) \in \mathcal{S} \times \mathcal{A}$, parameter averaging is equivalent to function averaging:
% \begin{equation}
%     Q_{\bar{\mathbf{M}}}(s,a) = \frac{1}{N}\sum_{i=1}^N Q_{\mathbf{M}^{(i)}}(s,a).
%     \label{eq:param_func_equiv}
% \end{equation}
% \end{theorem}

% \begin{proof}
% By definition of the HDC Q-function~\eqref{eq:hdc_q_function} and linearity of the inner product:
% \begin{align}
%     Q_{\bar{\mathbf{M}}}(s,a)
%     &= \langle \Phi(s,a), \bar{\mathbf{M}} \rangle \\
%     &= \left\langle \Phi(s,a), \frac{1}{N}\sum_{i=1}^N \mathbf{M}^{(i)} \right\rangle \\
%     &= \frac{1}{N}\sum_{i=1}^N \langle \Phi(s,a), \mathbf{M}^{(i)} \rangle \quad \text{(linearity of $\langle\cdot,\cdot\rangle$)} \\
%     &= \frac{1}{N}\sum_{i=1}^N Q_{\mathbf{M}^{(i)}}(s,a). \qquad \qed
% \end{align}
% \end{proof}

% This theorem justifies the direct parameter averaging strategy in Algorithm~\ref{alg:fedqhd_shared}
% and distinguishes HDC-based Q-learning from deep neural network approaches, where parameter
% averaging does not preserve function-level aggregation due to nonlinear activation functions.

% \subsection{Extension: Heterogeneous Encoders via Regression-Based Aggregation}
% \label{subsec:heterogeneous_extension}


\subsubsection{Regression-Based Aggregation Protocol}

The key insight is that while parameters cannot be directly averaged across different
feature spaces, Q-\textit{functions} can be averaged pointwise. By exploiting the linearity
of HDC Q-functions in parameter space, we can solve a regression problem to find global
parameters that approximate the averaged Q-function.

\textbf{Anchor-based aggregation:} The server maintains:
\begin{itemize}
    \item A global encoder $\Phi_{\text{global}}: \mathcal{S}\times\mathcal{A}\to\mathbb{C}^D$
    \item A set of anchor state-action pairs $\mathcal{X} = \{(s_j, a_j)\}_{j=1}^{L}$
    sampled from a representative distribution
\end{itemize}

\textbf{Upload phase:} Instead of uploading parameter vectors, each agent evaluates its
local Q-function on the anchor set:
\begin{equation}
    \mathbf{q}^{(i)}_k = \left[ Q_{\mathbf{M}^{(i)}_{k,K}}^{(i)}(s_j, a_j) \right]_{j=1}^{L}
    = \left[ \langle \Phi^{(i)}(s_j, a_j), \mathbf{M}^{(i)}_{k,K} \rangle \right]_{j=1}^{L}
    \in \mathbb{R}^L.
    \label{eq:q_value_upload_het}
\end{equation}
This requires $O(L)$ communication per agent, which can be more efficient than $O(D)$
parameter uploads when $L \ll D$.

\textbf{Aggregation phase:} The server computes the target averaged Q-values:
\begin{equation}
    \bar{\mathbf{q}}_k = \frac{1}{N}\sum_{i=1}^{N} \mathbf{q}^{(i)}_k,
    \label{eq:averaged_q_het}
\end{equation}
and solves a least-squares regression to find global parameters that best approximate
$\bar{\mathbf{q}}_k$ under $\Phi_{\text{global}}$:
\begin{equation}
    \mathbf{M}_{k+1} = \arg\min_{\mathbf{M} \in \mathbb{C}^D}
    \sum_{j=1}^{L} \left\|
    \langle \Phi_{\text{global}}(s_j, a_j), \mathbf{M} \rangle - \bar{q}_{k,j}
    \right\|^2.
    \label{eq:regression_aggregation_het}
\end{equation}

This has a closed-form solution via the Moore-Penrose pseudoinverse:
\begin{equation}
    \mathbf{M}_{k+1} = \mathbf{\Psi}^\dagger \bar{\mathbf{q}}_k,
    \label{eq:closed_form_aggregation_het}
\end{equation}
where $\mathbf{\Psi} = [\Phi_{\text{global}}(s_1,a_1), \dots, \Phi_{\text{global}}(s_L,a_L)]^T \in \mathbb{C}^{L \times D}$.

\subsubsection{Modified Algorithm for Heterogeneous Encoders}

Algorithm~\ref{alg:fedqhd_heterogeneous} presents the complete protocol with regression-based
aggregation.

\begin{algorithm}[t]
\caption{FedQHD: Heterogeneous Encoder Extension}
\label{alg:fedqhd_heterogeneous}
\begin{algorithmic}[1]
\STATE \textbf{Input:} Rounds $T$, local steps $K$, agents $N$, anchor size $L$, learning rates $\{\alpha_{k,u}\}$
\STATE \textbf{Server Initialize:}
\STATE \quad Global encoder $\Phi_{\text{global}}: \mathcal{S}\times\mathcal{A}\to\mathbb{C}^D$
\STATE \quad Global model $\mathbf{M}_0 = \mathbf{0} \in \mathbb{C}^D$
\STATE \quad Anchor set $\mathcal{X} = \{(s_j, a_j)\}_{j=1}^{L}$
\STATE \quad Feature matrix $\mathbf{\Psi} = [\Phi_{\text{global}}(s_1,a_1), \dots, \Phi_{\text{global}}(s_L,a_L)]^T \in \mathbb{C}^{L \times D}$
\STATE \textbf{Client Initialize:} Each agent $i$ initializes or adapts local encoder $\Phi^{(i)}$
\STATE
\FOR{round $k = 0, 1, \ldots, T-1$}
    \STATE \textbf{[Server Broadcast]} Send $(\mathbf{M}_k, \Phi_{\text{global}}, \mathcal{X})$ to all agents
    \STATE
    \STATE \textbf{[Client Local Learning]} (in parallel for all agents $i = 1, \ldots, N$)
    \STATE \quad Initialize: $\mathbf{M}^{(i)}_{k,0} = \mathbf{M}_k$
    \STATE \quad \textbf{(Optional) Adapt encoder:} Fine-tune $\Phi^{(i)}$ based on $\Phi_{\text{global}}$
    \FOR{local step $u = 0, 1, \ldots, K-1$}
        \STATE \quad Observe state $s_t^{(i)}$, select action $a_t^{(i)}$ via $\epsilon$-greedy
        \STATE \quad Execute $a_t^{(i)}$, observe $r_t^{(i)}, s_{t+1}^{(i)}$, select next action $a_{t+1}^{(i)}$
        \STATE \quad Encode: $\phi_t = \Phi^{(i)}(s_t^{(i)}, a_t^{(i)})$, $\phi_{t+1} = \Phi^{(i)}(s_{t+1}^{(i)}, a_{t+1}^{(i)})$
        \STATE \quad TD error: $\delta_t^{(i)} = r_t^{(i)} + \gamma \langle \phi_{t+1}, \mathbf{M}^{(i)}_{k,u} \rangle - \langle \phi_t, \mathbf{M}^{(i)}_{k,u} \rangle$
        \STATE \quad \textbf{Update:} $\mathbf{M}^{(i)}_{k,u+1} = \mathbf{M}^{(i)}_{k,u} + \alpha_{k,u} \delta_t^{(i)} \phi_t$
    \ENDFOR
    \STATE \quad \textbf{Evaluate on anchors:} $\mathbf{q}^{(i)}_k = [\langle \Phi^{(i)}(s_j, a_j), \mathbf{M}^{(i)}_{k,K} \rangle]_{j=1}^{L}$
    \STATE \quad \textbf{Upload:} Send Q-value vector $\mathbf{q}^{(i)}_k \in \mathbb{R}^L$ to server
    \STATE
    \STATE \textbf{[Server Aggregation]}
    \STATE \quad Receive $\{\mathbf{q}^{(i)}_k\}_{i=1}^{N}$ from all agents
    \STATE \quad \textbf{Average Q-values:} $\bar{\mathbf{q}}_k = \frac{1}{N}\sum_{i=1}^{N} \mathbf{q}^{(i)}_k$
    \STATE \quad \textbf{Regression:} $\mathbf{M}_{k+1} = \mathbf{\Psi}^\dagger \bar{\mathbf{q}}_k$
    \STATE \quad \textbf{(Optional)} Project: $\mathbf{M}_{k+1} \leftarrow \text{Proj}_{\mathcal{B}(G)}(\mathbf{M}_{k+1})$
\ENDFOR
\STATE
\STATE \textbf{Output:} Global Q-function $Q_{\mathbf{M}_T}(s,a) = \langle \Phi_{\text{global}}(s,a), \mathbf{M}_T \rangle$
\end{algorithmic}
\end{algorithm}

\subsubsection{Theoretical Guarantee}

\begin{theorem}[Regression-Based Aggregation Correctness]
\label{thm:regression_aggregation}
Let $\{\Phi^{(i)}\}_{i=1}^N$ be heterogeneous local encoders and $\{\mathbf{M}^{(i)}\}_{i=1}^N$ be
corresponding parameter vectors. Let $\Phi_{\text{global}}$ be the global encoder and
$\mathcal{X} = \{(s_j, a_j)\}_{j=1}^{L}$ be anchor points. Define the target averaged Q-function as
\begin{equation}
    \bar{Q}(s,a) = \frac{1}{N}\sum_{i=1}^N Q_{\mathbf{M}^{(i)}}^{(i)}(s,a),
\end{equation}
where $Q_{\mathbf{M}^{(i)}}^{(i)}(s,a) = \langle \Phi^{(i)}(s,a), \mathbf{M}^{(i)} \rangle$.
Let $\mathbf{M}_{\text{reg}}$ be the solution to~\eqref{eq:regression_aggregation_het}.
Then, $\mathbf{M}_{\text{reg}}$ minimizes the mean squared error on the anchor set:
\begin{equation}
    \mathbf{M}_{\text{reg}} = \arg\min_{\mathbf{M}} \sum_{j=1}^{L}
    \left| \langle \Phi_{\text{global}}(s_j, a_j), \mathbf{M} \rangle - \bar{Q}(s_j, a_j) \right|^2.
\end{equation}
Moreover, if $L \geq D$ and the anchor features span $\mathbb{C}^D$, then $Q_{\mathbf{M}_{\text{reg}}}$
exactly interpolates $\bar{Q}$ on $\mathcal{X}$.
\end{theorem}

\begin{proof}
The regression objective is a standard least-squares problem with linear model
$\langle \Phi_{\text{global}}(s_j, a_j), \mathbf{M} \rangle$. The closed-form solution
follows from normal equations. When $L \geq D$ and $\text{rank}(\mathbf{\Psi}) = D$,
the pseudoinverse satisfies $\mathbf{\Psi}^\dagger \mathbf{\Psi} = \mathbf{I}_D$, yielding
exact interpolation on $\mathcal{X}$. \qed
\end{proof}

\textbf{Remark:} Theorem~\ref{thm:regression_aggregation} shows that regression-based
aggregation preserves the averaged Q-function on anchor points, enabling heterogeneous
 model aggregation. The choice of anchor distribution and size $L$ controls the
approximation quality and generalization to unseen states. In practice, anchors can be
sampled from a replay buffer or from the empirical state-action distribution observed
during training.

